\documentclass[12pt]{article}
\usepackage[margin=0.9in]{geometry}
\usepackage[utf8]{inputenc}
\usepackage[T1]{fontenc}
\usepackage{array}
\usepackage{tabularx}
\usepackage{xltabular}
\usepackage{booktabs}
\usepackage[dvipsnames]{xcolor}
\usepackage{hyperref}
\usepackage{enumitem}
\usepackage{parskip}
\usepackage{setspace}
\onehalfspacing
\definecolor{navyBlue}{HTML}{003087}
\definecolor{brickRed}{HTML}{B22222}
\definecolor{tableHeader}{HTML}{003087}
\definecolor{altRow}{HTML}{F0F4FF}
\newcolumntype{P}[1]{>{\raggedright\arraybackslash}p{#1}}
\newcolumntype{L}{>{\raggedright\arraybackslash}X}
\renewcommand{\arraystretch}{1.16}
\hypersetup{colorlinks=true,linkcolor=navyBlue,urlcolor=brickRed}
\setlist[itemize]{leftmargin=*,itemsep=2pt,topsep=3pt}

\begin{document}

\begin{center}
{\Large\bfseries\color{navyBlue} ED452B Part 2 Level 4 Final Addendum}\par
{\large Advocacy, Professional Growth, Collaboration, and Research Strengthening}\par
{\normalsize Minecraft Education Coding FUNdamentals \textperiodcentered\ Piter Garcia \textperiodcentered\ May 2, 2026}\par
\end{center}

\section*{Purpose}

This final addendum addresses the last two rubric risk areas identified in the final review of the Part 2 Revised Innovative Unit Plan: \textbf{Unit Enactment Commentary} and \textbf{CEC Collaboration}. It also adds a concise private-data note and two peer-reviewed research connections to strengthen the theoretical framework.

\section{Private Data and Evidence Note for Submission}

The public portfolio uses synthesized and de-identified evidence only. Raw TeachingPlacement transcripts, student names, IEP documents, and identifiable work samples are intentionally excluded from the public repository. Fuller evidence records exist in private course and placement materials and were used only to identify patterns of access, planning, debugging, collaboration, and explanation. This distinction protects student privacy while still allowing the unit plan to show how data informed instructional decisions.

\section{Section 9 Addendum: Advocacy and Professional Growth}

The following language is written to be inserted at the end of Section 9, \emph{Unit Enactment Commentary and Critical Evaluation}.

\begin{quote}
A future enactment of this unit also requires advocacy. Minecraft Education can be misunderstood as a reward, enrichment activity, or behavior-management tool if the teacher does not clearly name the computer science and inclusion goals. I would advocate with administrators, colleagues, and families by explaining that the unit is designed to make algorithmic thinking visible while also documenting the access conditions students need to participate. The central message would be that inclusive technology integration is not lowering expectations; it is clarifying what the task is actually measuring. I would share the assessment categories --- startup access, task-goal understanding, planning, code accuracy, debugging, explanation, and support use --- so colleagues can see how the unit protects rigor while expanding access.

To keep building capacity, I would seek out professional resources from the CSTA community, Minecraft Education educator materials, Warner faculty feedback, and peer-reviewed research on inclusive computational thinking. I would also use post-lesson evidence meetings with the cooperating teacher to review which supports built independence and which supports should be faded. This advocacy matters because innovative practices only become sustainable when they are explainable to other professionals, connected to standards, and supported by evidence of student learning.
\end{quote}

\section{Section 10 Addendum: Cooperating Teacher Debrief Example}

The following language is written to be inserted in Section 10, \emph{Collaboration}.

\begin{quote}
A concrete collaboration point occurred after the debugging-focused portion of the Minecraft sequence. In debriefing the lesson, the cooperating teacher and I focused less on whether students finished the challenge and more on what kind of help made their thinking visible. We discussed that some students needed a precise prompt, such as counting spaces, restating the goal, or identifying what they expected the Agent to do before changing the code. That conversation helped clarify the difference between productive scaffolding and taking over the task. The next revision therefore strengthened the hints-before-rescue routine, the one-change debugging protocol, and the expectation that students explain the bug or fix after receiving support.

This debrief also created an opportunity for improvement. We noticed that students with strong Minecraft familiarity could become powerful peer supports, but only if the helper role was structured. Without structure, peer help could turn into one student controlling another student's keyboard. The final unit therefore defines the Expert Helper and Navigator roles more carefully: helpers ask questions, point to resources, and support planning, but they do not complete the code for another student. This collaboration shaped both the instructional routine and the assessment system because it made role participation, independence, and support use visible evidence of learning.
\end{quote}

\section{Research Strengthening for the Theoretical Framework}

The following paragraph can be inserted into Section 3, \emph{Theoretical Framework}, after the asset-based CS pedagogy subsection.

\begin{quote}
The unit is also supported by peer-reviewed research on inclusive and culturally responsive computer science education. Israel, Pearson, Tapia, Wherfel, and Reese's cross-case analysis of school-wide computational thinking found that computing integration can support diverse and struggling learners when teachers receive support, take instructional risks, and design collaborative learning structures. Scott, Sheridan, and Clark's theory of culturally responsive computing reinforces the need to treat students' identities, communities, and digital practices as resources for computing participation. These studies support the unit's design choices: Minecraft is used as an asset-based entry point, collaboration is structured rather than assumed, and assessment is distributed across code artifacts, explanation, debugging, and support use.
\end{quote}

\section{Level 4 Rubric Closure Matrix}

\begin{xltabular}{\textwidth}{|P{1.6in}|P{2.2in}|L|}
\hline
\rowcolor{tableHeader}\color{white}\bfseries Rubric Concern & \color{white}\bfseries Final Addition & \color{white}\bfseries Why This Moves Toward Level 4 \\
\hline
Unit Enactment Commentary & Advocacy paragraph explaining how to defend Minecraft as rigorous inclusive CS instruction. & Addresses the Level 4 expectation that the teacher can advocate for innovative practices on students' behalf. \\
\hline
\rowcolor{altRow} Professional Growth & Identifies CSTA, Minecraft Education, Warner feedback, research, and CT debriefs as capacity-building resources. & Shows how future practice will be strengthened through professional resources rather than isolated reflection. \\
\hline
CEC Collaboration & Adds a concrete post-lesson debrief example with the cooperating teacher. & Shows discussion of lessons learned, successes, and opportunities for improvement. \\
\hline
\rowcolor{altRow} Student Data Privacy & Clarifies that public evidence is synthesized and fuller records remain private. & Protects students while explaining why public artifacts use aggregate and de-identified records. \\
\hline
Research/Theory & Adds Israel et al. and Scott et al. as peer-reviewed sources. & Strengthens the already solid Karten/UDL/CEC framework with inclusive CS and culturally responsive computing scholarship. \\
\hline
\end{xltabular}

\section{References}

Israel, M., Pearson, J. N., Tapia, T., Wherfel, Q. M., \& Reese, G. (2015). Supporting all learners in school-wide computational thinking: A cross-case qualitative analysis. \emph{Computers \& Education, 82}, 263--279. \url{https://doi.org/10.1016/j.compedu.2014.11.022}

Scott, K. A., Sheridan, K. M., \& Clark, K. (2015). Culturally responsive computing: A theory revisited. \emph{Learning, Media and Technology, 40}(4), 412--436. \url{https://doi.org/10.1080/17439884.2014.924966}

\end{document}
